% Shared exercise chunk: l2-threshold-sampling
\begin{exercisechunk}
\item \textbf{What do the sampling assumptions buy?}
\label[exercise]{ex:threshold-sampling}
Fix a threshold target $f_k$. Let $\cA$ be the learner that returns
$f_{\kest}$, where $\kest$ is the smallest observed positive index,
or $N+1$ if none is observed.
\begin{enumerate}[(a)]
\item \emph{Arbitrary sampling.} Allow any joint distribution of the training
and test indices, as in the original boxes problem. Prove
\[
R(\cA,f_k)=\P\bigl(X\ge k,\;
X_i\notin\{k,\ldots,X\}\text{ for every }i\bigr).
\]
Prove this by showing equality of the corresponding events, before taking
probabilities.

\item \emph{Dependent training data with guaranteed coverage.}
Now suppose the test index is an independent draw from a distribution $P$.
The training indices may be dependent. Assume that, for some
$\alpha\in(0,1]$, every interval $I\subseteq[N]$ satisfies
\[
\P(X_i\in I\mid X_1,\ldots,X_{i-1})\ge\alpha P(I)
\qquad\text{almost surely}
\]
for each $i$, with unconditional probability when $i=1$.
Let $L=P(\{k,\ldots,\kest-1\})$ be the learned predictor's error mass,
where the set is empty if $\kest=k$.
Adapt the proof of \cref{prop:threshold} to show, for $\epsilon\in(0,1)$,
\[
\P(L>\epsilon)\le(1-\alpha\epsilon)^m\le e^{-\alpha m\epsilon}.
\]
\emph{Hint: bound the probability of missing a fixed interval at every step
by conditioning successively on the previous indices.}

\item \emph{Interpret the assumption.} \begin{enumerate}[label=(\roman*),leftmargin=*,beginpenalty=10000]
\item Verify the conditional coverage
inequality for independent draws from $P$ with $\alpha=1$.
\item Construct a
sampling procedure whose training indices are not independent and verify
the inequality for some $\alpha\in(0,1)$.
\item What sample size suffices to
ensure $L\le\epsilon$ with probability at least $1-\delta$, for
$\delta\in(0,1)$?
\end{enumerate}

\item \emph{Matching marginal distributions is insufficient.}
Let $N=2$, $f=f_1$, and $X_1=\cdots=X_m=Z$ for $m\ge1$, where $Z$ is
uniform on $\{1,2\}$. Let the test index be independently uniform.
Compute the learner's risk for every $m\ge1$. Use this expression to
determine whether increasing $m$ reduces the risk in this example.
\end{enumerate}

Independence of the test input makes its distribution $P$ the appropriate
measure of error after observing the training data. The coverage condition
ensures that each additional observation has a chance of revealing a
troublesome interval, even when the training observations are dependent.
Independent sampling from $P$ is one way to obtain this guarantee.
\end{exercisechunk}
